% ===================== CHAPITRE 1 : CONTEXTE GÉNÉRAL =====================

\section*{Introduction}
Ce chapitre situe le projet~: présentation de l'organisme d'accueil, analyse de
l'existant, formulation de la problématique et des objectifs, puis cahier des
charges.

\section{Présentation de l'organisme d'accueil}
\label{sec:organisme}

\begin{tcolorbox}[colback=codebg,colframe=subtitle,title=À compléter]
Section à rédiger avec les informations officielles de \rapportEntreprise\
(dénomination, forme juridique, création, effectif, activité, organigramme).
Le tableau ci-dessous sert de trame.
\end{tcolorbox}

\begin{table}[H]
\centering
\renewcommand{\arraystretch}{1.25}
\begin{tabularx}{\linewidth}{>{\bfseries}p{4cm} X}
\toprule
Raison sociale     & \rapportEntreprise \\
Forme juridique    & [SARL / SUARL / SA] \\
Date de création   & [\ldots] \\
Siège social       & [Adresse, gouvernorat] \\
Activité           & Édition de logiciels / services numériques \\
Effectif           & [\ldots] \\
\bottomrule
\end{tabularx}
\caption{Fiche d'identité de l'organisme d'accueil}
\label{tab:fiche-entreprise}
\end{table}

\section{Étude de l'existant}
\label{sec:existant}

La diffusion d'annonces immobilières en Tunisie repose sur des portails de
petites annonces généralistes, des sites d'agences et des réseaux sociaux.
L'analyse de ces solutions fait ressortir les limites suivantes~:

\begin{table}[H]
\centering
\renewcommand{\arraystretch}{1.25}
\begin{tabularx}{\linewidth}{>{\bfseries}p{4.2cm} X}
\toprule
Localisation approximative & Champ « ville » ou « quartier » en texte libre, sans position sur carte. \\
Recherche peu ergonomique  & Filtres limités, pas de recherche « autour de moi » ni par zone. \\
Découpage territorial ignoré & Le triplet gouvernorat / délégation / localité n'est pas exploité de façon structurée. \\
Qualité des annonces        & Doublons, annonces obsolètes, peu de modération. \\
Expérience mobile           & Interfaces non pensées « mobile-first ». \\
\bottomrule
\end{tabularx}
\caption{Limites des solutions existantes}
\label{tab:limites-existant}
\end{table}

\section{Problématique et objectifs}
\label{sec:problematique}

Il s'agit de concevoir une plateforme permettant de \emph{publier, géolocaliser
et rechercher} des biens immobiliers en Tunisie, en s'appuyant sur le découpage
territorial national et sur une carte interactive. Les objectifs sont~:

\begin{itemize}[label=\textbullet,itemsep=1pt]
  \item offrir une \textbf{recherche cartographique} fluide (déplacement, zoom,
        regroupement de marqueurs, recherche par zone)~;
  \item permettre la \textbf{publication géolocalisée} d'une annonce par
        placement d'un point sur la carte et géocodage inverse de l'adresse~;
  \item structurer les données autour du référentiel
        \textbf{gouvernorat $\rightarrow$ délégation $\rightarrow$ localité}~;
  \item gérer plusieurs \textbf{profils d'utilisateurs} et leurs droits~;
  \item proposer une expérience \textbf{responsive}, pensée pour le mobile.
\end{itemize}

\section{Cahier des charges}
\label{sec:cahier-des-charges}

\begin{table}[H]
\centering
\renewcommand{\arraystretch}{1.2}
\begin{tabularx}{\linewidth}{>{\bfseries}p{1.3cm} p{3.2cm} X}
\toprule
Réf. & Module & Description \\
\midrule
BF-01 & Comptes         & Inscription, connexion, profil, rôles. \\
BF-02 & Annonces        & Création, modification, suppression, publication. \\
BF-03 & Géolocalisation & Placement sur carte, latitude/longitude, géocodage inverse. \\
BF-04 & Référentiel géo. & Sélection gouvernorat / délégation / localité, recherche de lieu. \\
BF-05 & Recherche       & Filtres (catégorie, type, prix, surface, pièces\ldots), recherche sur carte. \\
BF-06 & Détail bien     & Fiche complète, galerie photos, carte de situation, contact. \\
BF-07 & Favoris \& alertes & Enregistrer un bien, sauvegarder une recherche, alerte de prix. \\
BF-08 & Comparateur     & Comparer plusieurs biens. \\
BF-09 & Espace agence   & Portefeuille d'annonces, agents, tableau de bord. \\
BF-10 & Modération      & Validation / refus des annonces, gestion des comptes (admin). \\
\bottomrule
\end{tabularx}
\caption{Besoins fonctionnels}
\label{tab:besoins-fonctionnels}
\end{table}

\begin{table}[H]
\centering
\renewcommand{\arraystretch}{1.2}
\begin{tabularx}{\linewidth}{>{\bfseries}p{3cm} X}
\toprule
Exigence & Description \\
\midrule
Ergonomie     & Interface claire, « mobile-first », conforme aux heuristiques de Nielsen. \\
Performance   & Affichage carte et résultats $<$ 2~s~; regroupement des marqueurs. \\
Portabilité   & Application web responsive, navigateurs récents. \\
Sécurité      & Authentification, contrôle d'accès par rôle. \\
Maintenabilité & Architecture en couches, modèle de données extensible. \\
Conformité    & Respect des principes du RGPD. \\
\bottomrule
\end{tabularx}
\caption{Besoins non fonctionnels}
\label{tab:besoins-non-fonctionnels}
\end{table}

\section*{Conclusion}
Le cadre du projet est posé~: contexte, existant, problématique et cahier des
charges. Le chapitre suivant présente la solution proposée.
